\documentclass[11pt,a4paper]{article}

\usepackage[margin=1in]{geometry}
\usepackage{hyperref}
\usepackage{enumitem}
\usepackage{booktabs}
\usepackage{tabularx}
\usepackage{xcolor}
\usepackage{mdframed}

\hypersetup{colorlinks=true, linkcolor=blue!70!black, urlcolor=blue!70!black}

\title{Our Indoor Navigation Approach\\[0.3em]
\large High-level overview for team planning and work assignment}
\author{Project Working Draft}
\date{\today}

\begin{document}

\maketitle
\begin{center}
\fbox{\parbox{0.9\textwidth}{
\textbf{HISTORICAL DOCUMENT} \\[4pt]
High-level overview from March 2026. Useful for seeing the early framing, but it
has been superseded by the later narrative documents and the current system guide. \\[4pt]
\textit{For the current system, see: \texttt{CLAUDE.md},
\texttt{final\_system\_overview.tex}}
}}
\end{center}
\vspace{1em}

\section{Goal}

Build a \textbf{reliable indoor navigation system} for the EarthRover that can
move through NYU indoor spaces using \textbf{image goals} rather than GPS.

\begin{mdframed}
\textbf{Core idea:} use a \textbf{topological graph} for global navigation,
\textbf{MBRA} for local image-goal control, and a \textbf{safety layer} to
prevent unsafe actions.
\end{mdframed}


\section{High-Level System}

Our system has five main parts:

\begin{enumerate}[nosep]
  \item \textbf{Perception} --- read camera frames and identify where the robot is.
  \item \textbf{Map / Graph} --- store indoor places as connected image nodes.
  \item \textbf{Planning} --- compute which node the robot should go to next.
  \item \textbf{Control} --- use MBRA to move toward the next subgoal image.
  \item \textbf{Safety} --- override unsafe commands before sending them.
\end{enumerate}

\subsection*{Data flow}

\begin{center}
\fbox{\parbox{0.93\textwidth}{
Camera stream $\rightarrow$ localization in graph $\rightarrow$ path to target
node $\rightarrow$ next subgoal image $\rightarrow$ MBRA local control
$\rightarrow$ safety check $\rightarrow$ robot command
}}
\end{center}


\section{Why This Approach}

\begin{itemize}[nosep]
  \item The competition is \textbf{image-goal based}, so a place-based graph is
        more natural than GPS-style navigation.
  \item MBRA is best used as a \textbf{short-horizon local controller}, not as
        the whole long-range navigation system.
  \item A separate planning layer makes the system easier to debug and more reliable.
  \item A separate safety layer reduces the risk of trusting the learned model too much.
\end{itemize}


\section{Module Breakdown}

\subsection{Perception}

\textbf{Purpose:} understand what the robot currently sees and localize it in
the indoor environment.

\textbf{Responsibilities:}
\begin{itemize}[nosep]
  \item capture and buffer front camera frames,
  \item extract visual descriptors for current frames,
  \item match current view to graph node images,
  \item provide current node estimate and confidence,
  \item support visual checkpoint verification.
\end{itemize}

\textbf{Outputs:}
\begin{itemize}[nosep]
  \item current node estimate,
  \item candidate nearby nodes,
  \item similarity scores / verification results.
\end{itemize}


\subsection{Planning}

\textbf{Purpose:} decide where the robot should go next at the place level.

\textbf{Responsibilities:}
\begin{itemize}[nosep]
  \item maintain the topological graph,
  \item map checkpoint images to target nodes,
  \item compute graph paths from current node to target node,
  \item choose the next local subgoal node on that path,
  \item handle node switching and simple recovery logic.
\end{itemize}

\textbf{Outputs:}
\begin{itemize}[nosep]
  \item target path,
  \item current subgoal image,
  \item path progress state.
\end{itemize}


\subsection{Control}

\textbf{Purpose:} move the robot from its current view to the next nearby
subgoal image.

\textbf{Responsibilities:}
\begin{itemize}[nosep]
  \item maintain recent image context,
  \item prepare MBRA inputs,
  \item run MBRA inference,
  \item convert MBRA outputs into rover commands,
  \item manage local command smoothing and rate limits.
\end{itemize}

\textbf{Outputs:}
\begin{itemize}[nosep]
  \item proposed linear/angular commands,
  \item controller state and diagnostics.
\end{itemize}


\subsection{Safety}

\textbf{Purpose:} prevent collisions and unsafe actions.

\textbf{Responsibilities:}
\begin{itemize}[nosep]
  \item estimate depth or clearance from the current frame,
  \item evaluate whether the MBRA action is safe,
  \item clip, redirect, or stop when needed,
  \item expose safety status to the rest of the system.
\end{itemize}


\subsection{Integration}

\textbf{Purpose:} connect everything into one runnable indoor system.

\textbf{Responsibilities:}
\begin{itemize}[nosep]
  \item EarthRover SDK interface,
  \item runtime orchestration between modules,
  \item configuration and logging,
  \item debugging tools,
  \item evaluation scripts and test runs.
\end{itemize}


\section{Suggested Team Workstreams}

\begin{tabularx}{\textwidth}{p{0.18\textwidth} p{0.35\textwidth} X}
\toprule
\textbf{Area} & \textbf{Main task} & \textbf{Expected deliverable} \\
\midrule
Perception & Visual localization and checkpoint matching & Current-node estimator and goal verification pipeline \\
Planning & Graph construction and graph search & Indoor topological graph format, path planner, subgoal selector \\
Control & MBRA runtime controller & MBRA wrapper and local goal-following controller on EarthRover \\
Safety & Runtime obstacle override & Depth/clearance-based safety module \\
Integration & End-to-end system glue & One script or service that runs the full pipeline \\
Testing & Evaluation and failure analysis & Repeatable test routes, metrics, debug reports \\
\bottomrule
\end{tabularx}


\section{Immediate Next Tasks}

\subsection*{Perception team}
\begin{itemize}[nosep]
  \item decide the first localization baseline,
  \item define how node images are stored,
  \item define the first goal verification method.
\end{itemize}

\subsection*{Planning team}
\begin{itemize}[nosep]
  \item define graph node and edge format,
  \item build a simple graph from one short indoor route,
  \item implement shortest-path search and subgoal selection.
\end{itemize}

\subsection*{Control team}
\begin{itemize}[nosep]
  \item load MBRA,
  \item buffer the last few images,
  \item feed a manually selected nearby goal image,
  \item generate rover commands from MBRA outputs.
\end{itemize}

\subsection*{Safety team}
\begin{itemize}[nosep]
  \item build a minimal runtime safety check,
  \item define stop / slow / override behavior.
\end{itemize}

\subsection*{Integration team}
\begin{itemize}[nosep]
  \item connect SDK + perception + planning + control + safety,
  \item make a minimal runnable indoor demo.
\end{itemize}


\section{First Milestone}

\begin{mdframed}
\textbf{Milestone 1:} prove that the robot can move toward a nearby indoor
subgoal image using MBRA while the rest of the system provides only minimal
support.
\end{mdframed}

This is the right first milestone because it validates the local controller
before we invest heavily in graph infrastructure.


\section{Second Milestone}

\begin{mdframed}
\textbf{Milestone 2:} connect graph localization and path planning so MBRA no
longer receives a manually selected goal image, but an automatically selected
subgoal node image.
\end{mdframed}


\section{What Success Looks Like}

We should consider the approach healthy if:
\begin{itemize}[nosep]
  \item the robot can localize itself in the graph,
  \item the planner can select the correct next node,
  \item MBRA can drive to that local node image,
  \item the safety layer prevents obvious collisions,
  \item and the full loop can repeat over multiple checkpoints.
\end{itemize}


\section{Short Summary}

\begin{mdframed}
\textbf{Our approach in one sentence:} use perception to localize in a
topological graph, planning to choose the next node, MBRA to drive toward that
node image, and safety to prevent bad local actions.
\end{mdframed}

That gives us a clear division of labor:
\begin{itemize}[nosep]
  \item \textbf{Perception} finds where we are,
  \item \textbf{Planning} decides where we go next,
  \item \textbf{Control} gets us there with MBRA,
  \item \textbf{Safety} keeps us from doing something stupid,
  \item \textbf{Integration} makes the whole system real.
\end{itemize}

\end{document}
